
%%%%%%%%%%%%%%%%%%%%%%%%%%%%%%%%%%%%%%%%%%%%%%%%%%%%%%%%%%%%%%%
\begin{frame}[t]
\frametitle{2.1.1 1入力可制御正準形への変換(演習)}
%-------------------------------------------------------------

\noindent
{\bf[演習]}\nopagebreak

次のシステムを可制御正準形に変換せよ．
\begin{align}
 \frac{d}{dt}
  \left(\begin{array}{c}
     {x}_1 \\
	 {x}_2 \\
	 {x}_3 
	\end{array}\right)
  &=
  \left(\begin{array}{ccc}
     0      & 2      & 3     \\
     1      & -1     & -1    \\
	 0      & 1      & 1     
	\end{array}\right)
  \left(\begin{array}{c}
     {x}_1 \\
	 {x}_2 \\
	 {x}_3 
	\end{array}\right)
  +
  \left(\begin{array}{c}
     1 \\
	 0 \\
	 0 
	\end{array}\right)
  u 
\nonumber \\
	y &= 
  \left(\begin{array}{ccc}
     0      & 1      & 1 
	\end{array}\right)
  \left(\begin{array}{c}
     {x}_1 \\
	 {x}_2 \\
	 {x}_3 
	\end{array}\right)
\nonumber 
\end{align}
\end{frame}


%%%%%%%%%%%%%%%%%%%%%%%%%%%%%%%%%%%%%%%%%%%%%%%%%%%%%%%%%%%%%%%
\begin{frame}[t]
\frametitle{2.1.1 1入力可制御正準形への変換(演習)}
%-------------------------------------------------------------

可制御性行列
\[
 V = \left(\begin{array}{ccc}
            B & A B & A^2 B
	   \end{array}\right)
   = \left(\begin{array}{ccc}
            1 & 0 & 2 \\
            0 & 1 & -1 \\
            0 & 0 & 1 
	   \end{array}\right)
\]
この逆行列の行ベクトルを以下で表す．
\[
 V^{-1} 
    = \left(\begin{array}{ccc}
            1 & 0 & -2 \\
            0 & 1 & 1 \\
            0 & 0 & 1 
	\end{array}\right)
    = \left(\begin{array}{c}
	     l_1^T \\
	     l_2^T \\
	     l_3^T 
	\end{array}\right)
\]
よって最下行は
\[
 l_3^T
    = \left(\begin{array}{ccc}
            0 & 0 & 1
	\end{array}\right)
\]

\end{frame}


%%%%%%%%%%%%%%%%%%%%%%%%%%%%%%%%%%%%%%%%%%%%%%%%%%%%%%%%%%%%%%%
\begin{frame}[t]
\frametitle{2.1.1 1入力可制御正準形への変換(演習)}
%-------------------------------------------------------------

座標変換
\begin{align*} 
	\bar{x}
	& =
	\left(\begin{array}{c}
	     \bar{x}_1 \\
	     \bar{x}_2 \\
	     \bar{x}_3 
	\end{array}\right)
	=
	\left(\begin{array}{c}
	     l_3^T  \\
	     l_3^T A  \\
	     l_3^T A^2 
	\end{array}\right)
	x
	=
	\left(\begin{array}{ccc}
	     0 & 0 & 1  \\
	     0 & 1 & 1  \\
	     1 & 0 & 0
	\end{array}\right)
	x
\nonumber \\ &
	= Wx
	= T^{-1}x
\end{align*}
よって座標変換行列は
\[
	T= 
	\left(\begin{array}{ccc}
	     0  & 0 & 1  \\
	     -1 & 1 & 0  \\
	     1  & 0 & 0
	\end{array}\right)
\]
\[
x = T \bar{x}
\]

\end{frame}


%%%%%%%%%%%%%%%%%%%%%%%%%%%%%%%%%%%%%%%%%%%%%%%%%%%%%%%%%%%%%%%
\begin{frame}[t]
\frametitle{2.1.1 1入力可制御正準形への変換(演習)}
%-------------------------------------------------------------

座標変換
\begin{align*}
	\bar{A}
	& = T^{-1} A T
	=
	\left(\begin{array}{ccc}
	     0 & 0 & 1  \\
	     0 & 1 & 1  \\
	     1 & 0 & 0
	\end{array}\right)
	\left(\begin{array}{ccc}
	     0      & 2      & 3     \\
	     1      & -1     & -1    \\
	     0      & 1      & 1     
	\end{array}\right)
	\left(\begin{array}{ccc}
	     0  & 0 & 1  \\
	     -1 & 1 & 0  \\
	     1  & 0 & 0
	\end{array}\right)
\\
	&=
	\left(\begin{array}{ccc}
	     0 & 1 & 1  \\
	     1 & 0 & 0  \\
	     0 & 2 & 3
	\end{array}\right)
	\left(\begin{array}{ccc}
	     0  & 0 & 1  \\
	     -1 & 1 & 0  \\
	     1  & 0 & 0
	\end{array}\right)
\\
	&=
	\left(\begin{array}{ccc}
	     0  & 1 & 0  \\
	     0  & 0 & 1  \\
	     1  & 2 & 0
	\end{array}\right)
\\
	\bar{B}
	& = T^{-1} B
	=
	\left(\begin{array}{ccc}
	     0 & 0 & 1  \\
	     0 & 1 & 1  \\
	     1 & 0 & 0
	\end{array}\right)
	\left(\begin{array}{c}
	     1 \\
	     0 \\
	     0 
	\end{array}\right)
	=
	\left(\begin{array}{c}
	     0 \\
	     0 \\
	     1 
	\end{array}\right)
\\
	\bar{C}
	& = C T
	=
	\left(\begin{array}{ccc}
	     0      & 1      & 1 
	\end{array}\right)
	\left(\begin{array}{ccc}
	     0  & 0 & 1  \\
	     -1 & 1 & 0  \\
	     1  & 0 & 0
	\end{array}\right)
	=
	\left(\begin{array}{ccc}
	     0      & 1      & 0 
	\end{array}\right)
\end{align*}

\end{frame}


%%%%%%%%%%%%%%%%%%%%%%%%%%%%%%%%%%%%%%%%%%%%%%%%%%%%%%%%%%%%%%%
\begin{frame}[t]
\frametitle{演習}
%-------------------------------------------------------------

\noindent
{\bf[演習]}\nopagebreak

次のシステムを考える
\begin{align}
 \frac{d}{dt}
  \left(\begin{array}{c}
     {x}_1 \\
	 {x}_2
	\end{array}\right)
  &=
  \left(\begin{array}{cc}
     0      & a  \\
     1      & 1
	\end{array}\right)
  \left(\begin{array}{c}
     {x}_1 \\
	 {x}_2
	\end{array}\right)
  +
  \left(\begin{array}{c}
     1 \\
	 1 
	\end{array}\right)
  u 
\nonumber \\
	y &= 
  \left(\begin{array}{cc}
     0      & 1   
	\end{array}\right)
  \left(\begin{array}{c}
     {x}_1 \\
	 {x}_2 
	\end{array}\right)
\nonumber 
\end{align}
(1) 伝達関数を求めよ．
\\[3mm]

(2) システムが不可制御になる$a$を求めよ．
\\[3mm]

(3) (2)の場合に伝達関数に極零相殺が起こっていることを確認せよ．
\\[3mm]

(4) (2)の場合のシステムの可安定性について論ぜよ．

\end{frame}



%%%%%%%%%%%%%%%%%%%%%%%%%%%%%%%%%%%%%%%%%%%%%%%%%%%%%%%%%%%%%%%
\begin{frame}[t]
\frametitle{演習}
%-------------------------------------------------------------
(1) 伝達関数は
\begin{align*}
	G(s)&= C (sI-A)^{-1} B 
	=
	\left(\begin{array}{cc}
 		    0      & 1   
		\end{array}\right)
	\left(\begin{array}{cc}
	     s      & -a  \\
	     -1      & s-1
		\end{array}\right)^{-1}
	\left(\begin{array}{c}
	     1 \\
	     1 
	\end{array}\right)
\\
	&=
	\left(\begin{array}{cc}
 		    0      & 1   
		\end{array}\right)
	\frac{
	\left(\begin{array}{cc}
	     s-1      & a  \\
	     1      & s
		\end{array}\right)
	}{s(s-1)-a}
	\left(\begin{array}{c}
	     1 \\
	     1 
	\end{array}\right)
\\
	&=
	\frac{
	\left(\begin{array}{cc}
 		    0      & 1   
		\end{array}\right)
	\left(\begin{array}{c}
	     s-1 + a  \\
	     1   +  s
		\end{array}\right)
	}{s^2-s-a}
\\
	&=\frac{s+1}{s^2-s-a}
\end{align*}

\end{frame}



%%%%%%%%%%%%%%%%%%%%%%%%%%%%%%%%%%%%%%%%%%%%%%%%%%%%%%%%%%%%%%%
\begin{frame}[t]
\frametitle{演習}
%-------------------------------------------------------------
(2) 可制御性行列は
\begin{align*}
	V&= (B,AB) 
	=
	\left[
		\left(\begin{array}{c}
		     1 \\
		     1 
		\end{array}\right)
	,
		\left(\begin{array}{c}
		     a \\
		     2 
		\end{array}\right)
	\right]
\end{align*}
　　よって$a=2$のとき不可制御
\\[3mm]

(3) $a=2$の時の伝達関数は
\begin{align*}
	G(s)&= \frac{s+1}{s^2-s-a}
	= \frac{s+1}{s^2-s-2}
	= \frac{s+1}{(s+1)(s-2)}
\end{align*}
\ \\

(4) 極零相殺されている極は$-1$なので可安定

\end{frame}
\end{document}
